\documentclass[../main.tex]{subfiles}

\begin{document}

\section{Building Vanilla Neural Network From Scratch}

For our toy examples, we will create two things. The first one is
technically not a neural network as it does not contain conventional
neurons and structures. The second one is an actual neural net that
is very simple. In this section, I will share the entire code first,
then provide explanations.

\subsection{Simple Stochastic Gradient Descent-Based Learning Engine}

What we will make now is a simple stochastic gradient descent-based
learning engine. This isn't technically a neural network, as you will
see why in the code, but I decided to include it here as making
similar examples really helped me understand backprop with code.

Our task with this learning engine is simple. We are going to predict
$e^x$ for limited range with concatenated Maclaurin Series.
If you recall from your calculus courses, the following equation holds
by Taylor series.
\[
e^x
= \sum_{n = 0}^\infty \frac{1}{n!} x^n
= 1 + x + \frac{x^2}{2!} + \frac{x^3}{3!} + \cdots
\]
Here, for limited $x$, we will be approximating $e^x$ with the
following polynomial where $b$ is the bias and $w_i$ are the
weights.
\[
e^x
\approx b + w_1x + w_2x^2 + w_3x^3 + w_4x^4 + w_5x^5 + w_6x^6
    + w_7x^7 + w_8x^8
\]
We will get back to this on why our $x$ should be limited, but let's
see our structure first.

We will be ``training'' the engine with $200$ examples from $-1.00$ to
$0.99$ inclusive. Then, we will be testing it with $50$ validations
from $1.00$ to $1.49$ inclusive. This backprop machine is stochastic
as it will adhere to the following structure.

\begin{figure}[H]
\centering
\begin{tikzpicture}[
    box/.style={rectangle, draw, thick,
                minimum height=0.8cm},
    arrow/.style={->, thick, >=stealth},
    scale=0.8
]

    \node[box, fill=CadetBlue!20] at (0, 0) (input)
        {\scriptsize Single Input};
    \node[box, fill=Lavender!30] at (4, 0) (mse)
        {\scriptsize MSE};
    \node[box, fill=YellowGreen!40] at (8, 0) (backprop)
        {\scriptsize Backpropagation};
    \node[box, fill=YellowOrange!30] at (12, 0) (update)
        {\scriptsize Update};

    \draw[arrow] (input) -- (mse);
    \draw[arrow] (mse) -- (backprop);
    \draw[arrow] (backprop) -- (update);
    \draw[arrow, bend left=20] (update) to (forward);
\end{tikzpicture}
\caption{Overview of the Simple Learning Engine}
\end{figure}

As you can see, we will be using MSE for our loss function, and we
do not need activation functions like sigmoid and tanh because we
don't need one in our simple example. Okay enough talking and here is
the entire code.

\subsubsection{Source Code}

Below is the source code for the Simple Stochastic Gradient
Descent-Based Learning Engine for predicting $e^x$.

\begin{jupyter}
\begin{lstlisting}[style=jin]
import math
import random
import matplotlib.pyplot as plt
%matplotlib inline
\end{lstlisting}
\begin{lstlisting}[style=jout]
\end{lstlisting}
\end{jupyter}

\begin{jupyter}
\begin{lstlisting}[style=jin]
train_x = [i/100 for i in range(-100, 100)]
train_y = [math.exp(i/100) for i in range(-100, 100)]
test_x = [i/100 for i in range(100, 150)]
test_y = [math.exp(i/100) for i in range(100, 150)]
\end{lstlisting}
\begin{lstlisting}[style=jout]
\end{lstlisting}
\end{jupyter}

\begin{jupyter}
\begin{lstlisting}[style=jin]
class SGDBLE:
    def __init__(self):
        random.seed(3141592) # So that we get the same result
        self.ws = [random.uniform(-1, 1) for _ in range(8)]
        self.b = random.random()
    
    def forward(self, x):
        self.ypred = sum(w * x**k for w, k in zip(self.ws,range(1, 9))) + self.b
        return self.ypred

    def loss(self, y, ypred):
        self.loss_val = 0.5 * (y - ypred)**2
        return self.loss_val

    def grads(self, x, y, ypred):
        base_der = ypred - y # i.e. dL / dypred
        self.grads_w = [base_der * x**i for i in range(1, 9)]
        self.grad_b = base_der
        
    def update(self, lr):
        for i in range(8):
            self.ws[i] -= lr * self.grads_w[i]
        self.b -= lr * self.grad_b

    def train(self, epochs, lr):
        losses = []
        for epoch in range(epochs):
            loss_total = 0
            for i in range(len(train_x)):
                ypred = self.forward(train_x[i])
                loss_total += self.loss(train_y[i], ypred)
                self.grads(train_x[i], train_y[i], ypred)
                self.update(lr)
            loss_avg = loss_total / len(train_x)
            losses.append(loss_avg)

            if (epoch+1) % 100 == 0:
                print(f"Epoch: {epoch+1}/{epochs}, Loss: {loss_avg}")
            
        return losses

    def test(self):
        loss_total = 0
        for i in range(0, len(test_x), 2):
            ypred1 = self.forward(test_x[i])
            loss_total += self.loss(test_y[i], ypred1)

            if i + 1 < len(test_x):
                ypred2 = self.forward(test_x[i + 1])
                loss_total += self.loss(test_y[i + 1], ypred2)
    
                print(f"Test {i+1}                      Test {i+2}")
                print(f"y:     {test_y[i]}      y:     {test_y[i+1]}")
                print(f"ypred: {ypred1}     ypred: {ypred2}")
                print()
            else:
                print(f"Test {i+1}")
                print(f"y:     {test_y[i]}")
                print(f"ypred: {ypred1}")
    
        loss_total /= len(test_x)

    def print_params(self):
        coeff = [1 / math.factorial(i+1) for i in range(8)]
        for i in range(8):
            print(f"w{i+1}: {self.ws[i]}, Expected: {coeff[i]}")
        print(f"b: {self.b}, Expected: 1")

model = SGDBLE()
\end{lstlisting}
\begin{lstlisting}[style=jout]
\end{lstlisting}
\end{jupyter}

\begin{jupyter}
\begin{lstlisting}[style=jin]
def graph_loss(losses):
    plt.plot(losses)
    plt.xlabel("Epoch")
    plt.ylabel("Loss")
    plt.show() 
    
def graph_pred():
    xs = [i / 10 for i in range(-100, 200)]
    true = [math.exp(x) for x in xs]
    pred = [model.forward(x) for x in xs]

    plt.plot(xs, true, label="e^x")
    plt.plot(xs, pred, label="Prediction")
    plt.yscale("log")
    plt.legend()
    plt.show()
\end{lstlisting}
\begin{lstlisting}[style=jout]
\end{lstlisting}
\end{jupyter}

\begin{jupyter}
\begin{lstlisting}[style=jin]
losses = model.train(1000, 0.005)
graph_loss(losses)
model.print_params()
graph_pred()
\end{lstlisting}
\begin{lstlisting}[style=jout]
Epoch: 100/1000, Loss: 0.00013234985876534958
Epoch: 200/1000, Loss: 3.427320321762344e-05
Epoch: 300/1000, Loss: 2.4492707827970274e-05
Epoch: 400/1000, Loss: 2.2669862503550263e-05
Epoch: 500/1000, Loss: 2.172205317266906e-05
Epoch: 600/1000, Loss: 2.0984248749352728e-05
Epoch: 700/1000, Loss: 2.036538265870468e-05
Epoch: 800/1000, Loss: 1.9834658817269138e-05
Epoch: 900/1000, Loss: 1.937270927247297e-05
Epoch: 1000/1000, Loss: 1.8965386261105253e-05
\end{lstlisting}
\vspace{-\baselineskip}
\includegraphics[scale=0.65]{CS0104.26017-1}
\begin{lstlisting}[style=jout]
w1: 0.9444454598562942, Expected: 1.0
w2: 0.49986306356857685, Expected: 0.5
w3: 0.620431907367758, Expected: 0.16666666666666666
w4: 0.21618569290559692, Expected: 0.041666666666666664
w5: -0.9321256839917001, Expected: 0.008333333333333333
w6: -0.48709437200615985, Expected: 0.001388888888888889
w7: 0.5594822789317027, Expected: 0.0001984126984126984
w8: 0.33002989302732855, Expected: 2.48015873015873e-05
b: 0.9983077779493585, Expected: 1
\end{lstlisting}
\vspace{-\baselineskip}
\includegraphics[scale=0.65]{CS0104.26017-2}
\end{jupyter}

\begin{jupyter}
\begin{lstlisting}[style=jin]
model.test()
\end{lstlisting}
\begin{lstlisting}[style=jout]
Test 1                      Test 2
y:     2.718281828459045      y:     2.7456010150169163
ypred: 2.7495260176087557     ypred: 2.7867851769675642

Test 3                      Test 4
y:     2.7731947639642978      y:     2.801065834699079
ypred: 2.8257757734623574     ypred: 2.8666322590606814

Test 5                      Test 6
y:     2.82921701435156      y:     2.857651118063164
ypred: 2.909497162739583     ypred: 2.954521430346852

Test 7                      Test 8
y:     2.8863709892679585      y:     2.915379499976997
ypred: 3.001864774074873     ypred: 3.0516960317129964

Test 9                      Test 10
y:     2.944679551065524      y:     2.9742740725630656
ypred: 3.1041935358457007     ypred: 3.159545493165112

Test 11                      Test 12
y:     3.0041660239464334      y:     3.034358394435676
ypred: 3.2179503740678106     ypred: 3.2796173127071624

Test 13                      Test 14
y:     3.0648542032930024      y:     3.095656500124711
ypred: 3.3447665176737535     ypred: 3.4136296934778425

Test 15                      Test 16
y:     3.1267683651861553      y:     3.158192909689767
ypred: 3.4864504730090555     ypred: 3.5634848611499024

Test 17                      Test 18
y:     3.1899332761161845      y:     3.2219926385284996
ypred: 3.6450016897210107     ypred: 3.7312830839373126

Test 19                      Test 20
y:     3.2543742028896707      y:     3.287081207383118
ypred: 3.8226249405557278     ypred: 3.9193374178962683

Test 21                      Test 22
y:     3.3201169227365472      y:     3.3534846525490236
ypred: 4.02174543791974     ypred: 4.130189200546641

Test 23                      Test 24
y:     3.3871877336213347      y:     3.4212295362896734
ypred: 4.245024710403111     ypred: 4.366624316181152

Test 25                      Test 26
y:     3.4556134647626755      y:     3.4903429574618414
ypred: 4.495377262801668     ypred: 4.631690256570204

Test 27                      Test 28
y:     3.5254214873653824      y:     3.5608525623555205
ypred: 4.775988043516579     ypred: 4.928714001110956

Test 29                      Test 30
y:     3.5966397255692817      y:     3.6327865557528094
ypred: 5.090330743550226     ypred: 5.261320740809877

Test 31                      Test 32
y:     3.6692966676192444      y:     3.706173712210199
ypred: 5.442186951657918     ypred: 5.633453470828671

Test 33                      Test 34
y:     3.7434213772608627      y:     3.781043387568781
ypred: 5.8356661905556715     ypred: 6.049393476664138

Test 35                      Test 36
y:     3.819043505366336      y:     3.8574255306969745
ypred: 6.275226859424931     ypred: 6.513781739373128

Test 37                      Test 38
y:     3.896193301795215      y:     3.9353506954704733
ypred: 6.765698108295758     ypred: 7.0316412855945165

Test 39                      Test 40
y:     3.9749016274947477      y:     4.014850052994201
ypred: 7.312302670230667     ypred: 7.6084005084606225

Test 41                      Test 42
y:     4.0551999668446745      y:     4.095955404071176
ypred: 7.9206806775719905     ypred: 8.249917485831345

Test 43                      Test 44
y:     4.137120440251392      y:     4.178699191923246
ypred: 8.596914488856152     ypred: 8.962505322624688

Test 45                      Test 46
y:     4.220695816996552      y:     4.263114515168817
ypred: 9.347554553339133     ypred: 9.752958544358306

Test 47                      Test 48
y:     4.305959528345206      y:     4.349235141062741
ypred: 10.179646340417854     ypred: 10.628580569357089

Test 49                      Test 50
y:     4.392945680918757      y:     4.437095519003664
ypred: 11.100758361572902     ypred: 11.597212287422602
\end{lstlisting}
\end{jupyter}

\subsubsection{Detailed Explanation}

Now let's take a look at individual cells.

\begin{jupyter}
\begin{lstlisting}[style=jin]
import math
import random
import matplotlib.pyplot as plt
%matplotlib inline
\end{lstlisting}
\begin{lstlisting}[style=jout]
\end{lstlisting}
\end{jupyter}

Here, I tried to keep it as minimal as possible. Instead of
NumPy, I decided to use the math module for all math stuffs like
\verb|math.exp()|. We need \verb|random| as it would make our lives
easier than to manually set the eight weights and the bias.
Matplotlib of course since we want to see how it learns for limited
range of input.

\begin{jupyter}
\begin{lstlisting}[style=jin]
train_x = [i/100 for i in range(-100, 100)]
train_y = [math.exp(i/100) for i in range(-100, 100)]
test_x = [i/100 for i in range(100, 150)]
test_y = [math.exp(i/100) for i in range(100, 150)]
\end{lstlisting}
\begin{lstlisting}[style=jout]
\end{lstlisting}
\end{jupyter}

Setting up the dataset for our engine is also quite simple.
Instead of writing every single number, I just decided to loop
through them and store them in an array, and that is \verb|train_x|.
The values for \verb|train_y| is also quite straightforward as it
is $e^x$ for corresponding $x$. The same goes for the test set.
I tried to keep it $80\%$ and $20\%$ for training and validation,
but that is not necessary and it varies by models. However, I
decided to use that here because that's pretty much what most
people do when they first train. Now let's get into the fun
stuffs.

\begin{jupyter}
\begin{lstlisting}[style=jin]
class SGDBLE:
    def __init__(self):
        random.seed(3141592) # So that we get the same result
        self.ws = [random.uniform(-1, 1) for _ in range(8)]
        self.b = random.random()
\end{lstlisting}
\begin{lstlisting}[style=jout]
\end{lstlisting}
\end{jupyter}

Now I decided to store the core functions in the class SGDBLE
as we will be using \verb|self| to store attributes and objects.
Here we set the specific seed so that you will also get the same
result when you run the code in somewhere. Next we set weights
and biases and stored them in \verb|self| with our set random values.

\begin{jupyter}
\begin{lstlisting}[style=jin]
def forward(self, x):
    self.ypred = sum(w * x**k for w, k in zip(self.ws,range(1, 9))) + self.b
    return self.ypred

def loss(self, y, ypred):
    self.loss_val = 0.5 * (y - ypred)**2
    return self.loss_val
\end{lstlisting}
\begin{lstlisting}[style=jout]
\end{lstlisting}
\end{jupyter}

This is where we set some core functions. Forward is forward pass
and it is somewhat straight forward as that is how we defined it.
For the loss function, we are using MSE, i.e. $\frac{1}{2} (y -
\hat{y})^2$.

\begin{jupyter}
\begin{lstlisting}[style=jin]
def grads(self, x, y, ypred):
    base_der = ypred - y # i.e. dL / dypred
    self.grads_w = [base_der * x**i for i in range(1, 9)]
    self.grad_b = base_der

def update(self, lr):
    for i in range(8):
        self.ws[i] -= lr * self.grads_w[i]
    self.b -= lr * self.grad_b
\end{lstlisting}
\begin{lstlisting}[style=jout]
\end{lstlisting}
\end{jupyter}

This is the core part of our engine. Let's briefly look at the math
and write the loop out for some to see what's actually going on.
Consider the following equation by the chain rule.
\[
\frac{dL}{dw_i} = \frac{dL}{dy} \cdot \frac{dy}{dw_i}
\]
Notice that $\frac{dL}{dy} = \frac{d}{dy} (\frac{1}{2} (y -
\hat{y})^2) = y - \hat{y}$. This is \verb|base_der|. Continuing with
the chain rule, we can write the weights and the bias as the
following.
\begin{align*}
    \frac{dL}{dw_1}
    &= \frac{dL}{dy} \cdot \frac{dy}{dw_1}
    = \text{base\_der} \cdot x \\
    \frac{dL}{dw_2}
    &= \frac{dL}{dy} \cdot \frac{dy}{dw_2}
    = \text{base\_der} \cdot x^2 \\
    \frac{dL}{db}
    &= \frac{dL}{dy} \cdot \frac{dy}{db}
    = \text{base\_der} \cdot 1
\end{align*}
Now the update function. Recall $w_2 = w_1 - \eta \frac{dL}{dw_1}$
and $b_2 = b_1 - \eta \frac{dL}{db_1}$. That is exactly what we did,
but for weights we have to loop through each element.

\begin{jupyter}
\begin{lstlisting}[style=jin]
def train(self, epochs, lr):
    losses = []
    for epoch in range(epochs):
        loss_total = 0
        for i in range(len(train_x)):
            ypred = self.forward(train_x[i])
            loss_total += self.loss(train_y[i], ypred)
            self.grads(train_x[i], train_y[i], ypred)
            self.update(lr)
        loss_avg = loss_total / len(train_x)
        losses.append(loss_avg)

        if (epoch+1) % 100 == 0:
            print(f"Epoch: {epoch+1}/{epochs}, Loss: {loss_avg}")
        
    return losses
\end{lstlisting}
\begin{lstlisting}[style=jout]
\end{lstlisting}
\end{jupyter}

This is the training loop. Of course we don't want to stop at a
single epoch especially in stochastic setting. You can ignore
the \verb|losses| in the beginning as we only use them to graph.
For each epoch, we want to set the total loss to $0$ as they will
accumulate. Then, for all $x$ values in our training set, we will
predict, calculate the loss, backprop, an update as shown. Note
that we have to set \verb|epoch| and \verb|lr|, which is the
learning rate. Now we will take the average of the total loss and
print it every $100$ epoch to see how well our engine is learning.

\begin{jupyter}
\begin{lstlisting}[style=jin]
def test(self):
    loss_total = 0
    for i in range(0, len(test_x), 2):
        ypred1 = self.forward(test_x[i])
        loss_total += self.loss(test_y[i], ypred1)

        if i + 1 < len(test_x):
            ypred2 = self.forward(test_x[i + 1])
            loss_total += self.loss(test_y[i + 1], ypred2)

            print(f"Test {i+1}                      Test {i+2}")
            print(f"y:     {test_y[i]}      y:     {test_y[i+1]}")
            print(f"ypred: {ypred1}     ypred: {ypred2}")
            print()
        else:
            print(f"Test {i+1}")
            print(f"y:     {test_y[i]}")
            print(f"ypred: {ypred1}")

    loss_total /= len(test_x)

def print_params(self):
    coeff = [1 / math.factorial(i+1) for i in range(8)]
    for i in range(8):
        print(f"w{i+1}: {self.ws[i]}, Expected: {coeff[i]}")
    print(f"b: {self.b}, Expected: 1")

model = SGDBLE()
\end{lstlisting}
\begin{lstlisting}[style=jout]
\end{lstlisting}
\end{jupyter}

The test loop is also straight forward. It predicts, calculates the
loss, and print them as shown. Also, because we know the coefficients
for the Maclaurin Series, I decided to see if the parameters match
with the actual coefficients. Finally, the last line is to be able
use the model.

\begin{jupyter}
\begin{lstlisting}[style=jin]
def graph_loss(losses):
    plt.plot(losses)
    plt.xlabel("Epoch")
    plt.ylabel("Loss")
    plt.show()
    
def graph_pred():
    xs = [i / 10 for i in range(-100, 200)]
    true = [math.exp(x) for x in xs]
    pred = [model.forward(x) for x in xs]

    plt.plot(xs, true, label="e^x")
    plt.plot(xs, pred, label="Prediction")
    plt.yscale("log")
    plt.legend()
    plt.show()
\end{lstlisting}
\begin{lstlisting}[style=jout]
\end{lstlisting}
\end{jupyter}

This part is plotting with Matplotlib. Please refer to the second
note if you are confused! One thing though \verb|yscale| is necessary
as exponential function will get too large to show.

\begin{jupyter}
\begin{lstlisting}[style=jin]
losses = model.train(1000, 0.005)
graph_loss(losses)
model.print_params()
graph_pred()
\end{lstlisting}
\begin{lstlisting}[style=jout]
\end{lstlisting}
\end{jupyter}

This part is the actual training part. We set the epoch to $1000$
and learning rate to $0.005$. This learning rate might be too small
for the start, but I realized that this gives slightly better results
than $0.001$ in this setting. Then we graph and show loss, parameters,
and predictions.

\begin{jupyter}
\begin{lstlisting}[style=jin]
model.test()
\end{lstlisting}
\begin{lstlisting}[style=jout]
\end{lstlisting}
\end{jupyter}

This part is for testing. Now that we have seen the entire code and
explanation for it, let's now analyze our results.

\subsubsection{Analyzing the Results}

First, let's take a look at the first result.

\begin{jupyter}
\begin{lstlisting}[style=jin]
losses = model.train(1000, 0.005)
graph_loss(losses)
model.print_params()
graph_pred()
\end{lstlisting}
\begin{lstlisting}[style=jout]
Epoch: 100/1000, Loss: 0.00013234985876534958
Epoch: 200/1000, Loss: 3.427320321762344e-05
Epoch: 300/1000, Loss: 2.4492707827970274e-05
Epoch: 400/1000, Loss: 2.2669862503550263e-05
Epoch: 500/1000, Loss: 2.172205317266906e-05
Epoch: 600/1000, Loss: 2.0984248749352728e-05
Epoch: 700/1000, Loss: 2.036538265870468e-05
Epoch: 800/1000, Loss: 1.9834658817269138e-05
Epoch: 900/1000, Loss: 1.937270927247297e-05
Epoch: 1000/1000, Loss: 1.8965386261105253e-05
\end{lstlisting}
\vspace{-\baselineskip}
\includegraphics[scale=0.65]{CS0104.26017-1}
\begin{lstlisting}[style=jout]
w1: 0.9444454598562942, Expected: 1.0
w2: 0.49986306356857685, Expected: 0.5
w3: 0.620431907367758, Expected: 0.16666666666666666
w4: 0.21618569290559692, Expected: 0.041666666666666664
w5: -0.9321256839917001, Expected: 0.008333333333333333
w6: -0.48709437200615985, Expected: 0.001388888888888889
w7: 0.5594822789317027, Expected: 0.0001984126984126984
w8: 0.33002989302732855, Expected: 2.48015873015873e-05
b: 0.9983077779493585, Expected: 1
\end{lstlisting}
\vspace{-\baselineskip}
\includegraphics[scale=0.65]{CS0104.26017-2}
\end{jupyter}

This is what we got after the first training. The loss looks
pretty good actually. As you can see for each $100^\text{th}$ epoch,
the loss is decreasing well. From the graph, you can see the sudden
drop in losses; however, we can't necessarily say that this is a
grokking phenomena as it might just be model understanding the
pattern from underfitting.

Looking at the weights, some are good and some are bad. We will
get back to this later, but this is also one of the reason why
we cannot scale this engine. Similarly, we can see that our
polynomial fits very well for the small range, but not well as
the input increases.

\begin{jupyter}
\begin{lstlisting}[style=jin]
model.test()
\end{lstlisting}
\begin{lstlisting}[style=jout]
Test 1                      Test 2
y:     2.718281828459045      y:     2.7456010150169163
ypred: 2.7495260176087557     ypred: 2.7867851769675642

Test 3                      Test 4
y:     2.7731947639642978      y:     2.801065834699079
ypred: 2.8257757734623574     ypred: 2.8666322590606814

Test 5                      Test 6
y:     2.82921701435156      y:     2.857651118063164
ypred: 2.909497162739583     ypred: 2.954521430346852
...
Test 45                      Test 46
y:     4.220695816996552      y:     4.263114515168817
ypred: 9.347554553339133     ypred: 9.752958544358306

Test 47                      Test 48
y:     4.305959528345206      y:     4.349235141062741
ypred: 10.179646340417854     ypred: 10.628580569357089

Test 49                      Test 50
y:     4.392945680918757      y:     4.437095519003664
ypred: 11.100758361572902     ypred: 11.597212287422602
\end{lstlisting}
\end{jupyter}

For the training set, I only listed some of them as you can
refer to the previous pages. Notice that for first few tests,
the predictions are not bad. However, for Test 5, the engine
does not seem be effective and seems lost from Test 10.
We could actually run our training loop multiple times to see
what happens, and I will leave this to the readers. However,
the engine does not seem to learn well.

Here is the question. Is the engine overfitting right now? If you
said no, well... that is correct! Technically, we can see it as
overfitting, but what I think is more accurate is to say that this
is the limitation of the scalability of this engine and our method.
When predicting, we only used nine terms. From the parameters,
some were very close to the expectation and others were very bad.
This is because we are using a polynomial with finite term to
predict, whereas the real Maclaurin series uses infinite sum.
Therefore, it is techinically an overfitting to certain point in
a way that our engine tries to fit our polynomial in limited range
as shown in the graph, but this too could improve if we have the
computational ability to use infinite sum. Therefore, it may be
more accurate to say that it is the limitation of the engine and
scalability that causes high loss in testing.

We just completed our first model! Although it is very limited
in terms of scalability and usability, I hope this example was
very helpful in using backpropagation in code. For our next
example, we will build an actual simple neural net to predict
a nonlinear function.

\subsection{Vanilla Neural Network}

Now that we have created a simple backprop engine, let's increase
the number of hidden layers to predict $x_1 x_2$ for inputs $x_1$
and $x_2$. I thought it would be great to introduce nonlinearity
to our problem, and it works! Without further ado, let's get started
by taking a look at the source code.

\subsubsection{Source Code}

Below is the entire source code with the results.

\begin{jupyter}
\begin{lstlisting}[style=jin]
import math
import random
import matplotlib.pyplot as plt
%matplotlib inline
\end{lstlisting}
\begin{lstlisting}[style=jout]
\end{lstlisting}
\end{jupyter}

\begin{jupyter}
\begin{lstlisting}[style=jin]
# Predicting x1 * x2
def tanh(x):
    return (math.exp(x) - math.exp(-x)) / (math.exp(x) + math.exp(-x))

def tanh_der(x):
    return 1 - tanh(x)**2

def loss(y, ypred):
    return 0.5 * (y - ypred)**2


# Data
data_x = [[i/10, j/10] for i in range(-10, 11) for j in range(-10, 11)]
data_y = [(x1 * x2) for x1, x2 in data_x]

data = list(zip(data_x, data_y))
random.seed(3141592)
random.shuffle(data)

split = int(0.8 * len(data))
train_data = data[:split]
test_data = data[split:]

train_x = [x for x, y in train_data]
train_y = [y for x, y in train_data]
test_x = [x for x, y in test_data]
test_y = [y for x, y in test_data]
\end{lstlisting}
\begin{lstlisting}[style=jout]
\end{lstlisting}
\end{jupyter}

\begin{jupyter}
\begin{lstlisting}[style=jin]
class MLP:
    def __init__(self, layers):
        self.ws = []
        self.bs = []

        random.seed(3141592)
        for i in range(len(layers) - 1):
            layer_in = layers[i]
            layer_out = layers[i + 1]
            layer_w = [[random.uniform(-1, 1) for _ in range(layer_in)]
                       for _ in range(layer_out)]
            layer_b = [random.uniform(-1, 1) for _ in range(layer_out)]

            self.ws.append(layer_w)
            self.bs.append(layer_b)
            
    def forward(self, x):
        self.act = [x]
        self.zs = []
        
        # Layers
        for i in range(len(self.ws)):
            prev = self.act[-1]
            pre_layer = []
            act_layer = []

            # Neurons
            for j in range(len(self.ws[i])):
                pre = self.bs[i][j]
                for k in range(len(self.ws[i][j])):
                    pre += self.ws[i][j][k] * prev[k]
                
                pre_layer.append(pre)

                if i != len(self.ws) - 1:
                    act_layer.append(tanh(pre))
                else:
                    act_layer.append(pre)
                    
            self.zs.append(pre_layer)
            self.act.append(act_layer)
            
        return self.act[-1][0]

    def backprop(self, y):
        self.grad_w = [[[0.0 for _ in range(len(self.ws[i][j]))]
                        for j in range(len(self.ws[i]))]
                       for i in range(len(self.ws))]
        self.grad_b = [[0.0 for _ in range(len(self.bs[i]))]
                       for i in range(len(self.bs))]
        ypred = self.act[-1][0]
        delta = [0] * len(self.ws)
        delta[-1] = [ypred - y]

        # range(start, stop, step)
        for k in range(len(self.ws) - 2, -1, -1):
            delta[k] = []
            for i in range(len(self.ws[k])):
                dlda = 0.0
                for j in range(len(self.ws[k + 1])):
                    dlda += self.ws[k + 1][j][i] * delta[k + 1][j]
                delta[k].append(dlda * tanh_der(self.zs[k][i]))

        for i in range(len(self.ws)):
            prev = self.act[i]
            for j in range(len(self.ws[i])):
                self.grad_b[i][j] = delta[i][j]
                for k in range(len(prev)):
                    self.grad_w[i][j][k] = delta[i][j] * prev[k]

    def update(self, lr):
        for i in range(len(self.ws)):
            for j in range(len(self.ws[i])):
                for k in range(len(self.ws[i][j])):
                    self.ws[i][j][k] -= lr * self.grad_w[i][j][k]
                self.bs[i][j] -= lr * self.grad_b[i][j]

    def train(self, train_x, train_y, epochs, lr):
        losses = []
        for epoch in range(epochs):
            total_loss = 0

            for x, y in zip(train_x, train_y):
                ypred = self.forward(x)
                total_loss += loss(y, ypred)
                self.backprop(y)
                self.update(lr)

            average_loss = total_loss / len(train_x)
            losses.append(average_loss)
            
            if (epoch + 1) % 100 == 0:
                print(f"Epoch: {epoch+1}/{epochs}, Loss: {average_loss}")
                
        return losses


    def show_params(self):
        for i, (w, b) in enumerate(zip(self.ws, self.bs)):
            layer = "OUTPUT" if i == len(self.ws) - 1 else f"LAYER {i+1}"
            print(layer)
            for j, neuron_w in enumerate(w):
                weights_str = ", ".join(f"{x:.4f}" for x in neuron_w)
                print(f"  Neuron {j+1}:")
                print(f"    Weights: [{weights_str}]")
                print(f"    Bias = {b[j]:.4f}")
            print()

mlp = MLP([2, 8, 8, 1])
\end{lstlisting}
\begin{lstlisting}[style=jout]
\end{lstlisting}
\end{jupyter}

\begin{jupyter}
\begin{lstlisting}[style=jin]
losses = mlp.train(train_x, train_y, epochs=2000, lr=0.01)
print()

for x, y in zip(train_x[:10], train_y[:10]):
    print(x, y)
    print(f"     ypred: {mlp.forward(x)}")

plt.plot(losses)
plt.xlabel("epoch")
plt.ylabel("loss")
plt.show()

#mlp.show_params()
\end{lstlisting}
\begin{lstlisting}[style=jout]
Epoch: 100/2000, Loss: 0.0006795237745380497
Epoch: 200/2000, Loss: 0.0002780199905385479
Epoch: 300/2000, Loss: 0.00015780133554953304
Epoch: 400/2000, Loss: 9.996884725589866e-05
Epoch: 500/2000, Loss: 6.544552166502369e-05
Epoch: 600/2000, Loss: 4.373246439583505e-05
Epoch: 700/2000, Loss: 3.032175895325618e-05
Epoch: 800/2000, Loss: 2.2295750007135315e-05
Epoch: 900/2000, Loss: 1.7559747007674843e-05
Epoch: 1000/2000, Loss: 1.4723606896641237e-05
Epoch: 1100/2000, Loss: 1.2951440585334302e-05
Epoch: 1200/2000, Loss: 1.1773900112740833e-05
Epoch: 1300/2000, Loss: 1.0936643742787753e-05
Epoch: 1400/2000, Loss: 1.0302935457567813e-05
Epoch: 1500/2000, Loss: 9.79811628836463e-06
Epoch: 1600/2000, Loss: 9.379957580934361e-06
Epoch: 1700/2000, Loss: 9.023360655331247e-06
Epoch: 1800/2000, Loss: 8.712547584218995e-06
Epoch: 1900/2000, Loss: 8.437042172730164e-06
Epoch: 2000/2000, Loss: 8.189547717119376e-06

[1.0, -1.0] -1.0
     ypred: -0.9914314084230468
[0.8, -1.0] -0.8
     ypred: -0.8026975789984329
[0.8, -0.9] -0.7200000000000001
     ypred: -0.7229659139930072
[0.1, -0.8] -0.08000000000000002
     ypred: -0.07901899113208088
[0.3, -0.3] -0.09
     ypred: -0.08931568886644314
[-0.7, -0.1] 0.06999999999999999
     ypred: 0.07003157242570901
[-0.6, -0.5] 0.3
     ypred: 0.29863205953269367
[0.7, 0.8] 0.5599999999999999
     ypred: 0.5613061391881997
[-0.5, -0.4] 0.2
     ypred: 0.19629911113434118
[0.9, 0.8] 0.7200000000000001
     ypred: 0.7218073084017129
\end{lstlisting}
\vspace{-\baselineskip}
\includegraphics[scale=0.65]{CS0104.26017-3}
\end{jupyter}

\begin{jupyter}
\begin{lstlisting}[style=jin]
test_loss = 0
for x, y in zip(test_x, test_y):
    ypred = mlp.forward(x)
    test_loss += loss(y, ypred)

for x, y in zip(test_x[:10], test_y[:10]):
    ypred = mlp.forward(x)
    print(x, y)
    print(f"    ypred: {ypred}")

print("Test Loss:", test_loss / len(test_x))
\end{lstlisting}
\begin{lstlisting}[style=jout]
[-0.9, -0.3] 0.27
    ypred: 0.27224246623233583
[0.9, 1.0] 0.9
    ypred: 0.889162612694173
[0.1, 0.5] 0.05
    ypred: 0.050824340568330406
[1.0, 1.0] 1.0
    ypred: 0.9737260427905209
[0.6, -0.8] -0.48
    ypred: -0.4811987339447982
[0.6, -0.3] -0.18
    ypred: -0.17587974783858806
[0.6, 0.4] 0.24
    ypred: 0.23577905969192914
[0.5, -0.8] -0.4
    ypred: -0.4007379739663608
[0.5, 0.2] 0.1
    ypred: 0.0970149401772119
[0.9, -0.2] -0.18000000000000002
    ypred: -0.18071247280108016
Test Loss: 1.1432408362631248e-05
\end{lstlisting}
\end{jupyter}

\subsubsection{Detailed Explanation}

Now let's take a look at each block of the code to see their role
in our vanilla neural net.

\begin{jupyter}
\begin{lstlisting}[style=jin]
import math
import random
import matplotlib.pyplot as plt
%matplotlib inline
\end{lstlisting}
\begin{lstlisting}[style=jout]
\end{lstlisting}
\end{jupyter}

Here, I imported essential modules and libraries to do math, get
a randomized value, and plot the loss.

\begin{jupyter}
\begin{lstlisting}[style=jin]
# Predicting x1 * x2
def tanh(x):
    return (math.exp(x) - math.exp(-x)) / (math.exp(x) + math.exp(-x))

def tanh_der(x):
    return 1 - tanh(x)**2

def loss(y, ypred):
    return 0.5 * (y - ypred)**2
\end{lstlisting}
\begin{lstlisting}[style=jout]
\end{lstlisting}
\end{jupyter}

For this problem, I decided to use tanh for our activation function
as sigmoid squishes the values to 0 and 1, which is not ideal for
regression problem like ours and ReLU was not effective when I tried.
Since $\tanh(x) = \frac{e^x - e^{-x}}{e^x + e^{-x}}$ and $\tanh'(x) =
1 - \tanh^2(x)$, we did the same. For our loss function, I decided to
use MSE for simplicity again.

\begin{jupyter}
\begin{lstlisting}[style=jin]
# Data
data_x = [[i/10, j/10] for i in range(-10, 11) for j in range(-10, 11)]
data_y = [(x1 * x2) for x1, x2 in data_x]

data = list(zip(data_x, data_y))
random.seed(3141592)
random.shuffle(data)

split = int(0.8 * len(data))
train_data = data[:split]
test_data = data[split:]

train_x = [x for x, y in train_data]
train_y = [y for x, y in train_data]
test_x = [x for x, y in test_data]
test_y = [y for x, y in test_data]
\end{lstlisting}
\begin{lstlisting}[style=jout]
\end{lstlisting}
\end{jupyter}

This part is for making our datasets. I decided to divide $i$ and $j$
by $10$ because if I don't, math overflow error will be resulted
due to $\tanh$ function.

This time, I decided to shuffle data in a way that we will get the
same result when you try running this, and split the data into
$80\%$ for training and $20\%$ for validating.

\begin{jupyter}
\begin{lstlisting}[style=jin]
class MLP:
    def __init__(self, layers):
        self.ws = []
        self.bs = []

        random.seed(3141592)
        for i in range(len(layers) - 1):
            layer_in = layers[i]
            layer_out = layers[i + 1]
            layer_w = [[random.uniform(-1, 1) for _ in range(layer_in)]
                       for _ in range(layer_out)]
            layer_b = [random.uniform(-1, 1) for _ in range(layer_out)]

            self.ws.append(layer_w)
            self.bs.append(layer_b)
\end{lstlisting}
\begin{lstlisting}[style=jout]
\end{lstlisting}
\end{jupyter}

Again for our example, I decided to store our core functions in
a class since we want to reuse and store our attributes.
When this class is first called, we will create an empty array
and for \verb|ws| and \verb|bs| and store them in \verb|self|. Then,
for the randomized value, we will loop through each layers to store
a randomized value between $-1$ and $1$ for weights and biases.
Notice that unlike the bias which only loop through the neurons
in the output layer, we have to loop through the previous layer
too for weights since we have different weights for each previous
neuron for each output neuron. We will then append, or store, our
layer to \verb|ws| and \verb|bs|.

\begin{jupyter}
\begin{lstlisting}[style=jin]
def forward(self, x):
    self.act = [x]
    self.zs = []
    
    # Layers
    for i in range(len(self.ws)):
        prev = self.act[-1]
        pre_layer = []
        act_layer = []

        # Neurons
        for j in range(len(self.ws[i])):
            pre = self.bs[i][j]
            for k in range(len(self.ws[i][j])):
                pre += self.ws[i][j][k] * prev[k]
            
            pre_layer.append(pre)

            if i != len(self.ws) - 1:
                act_layer.append(tanh(pre))
            else:
                act_layer.append(pre)
                
        self.zs.append(pre_layer)
        self.act.append(act_layer)
        
    return self.act[-1][0]
\end{lstlisting}
\begin{lstlisting}[style=jout]
\end{lstlisting}
\end{jupyter}

This is the forward function. Here, I decided to create both
\verb|self.act| and \verb|self.zs| since we need both of them
for backprop function. As the name suggests, \verb|self.act| is the
value after $\tanh$ is applied to \verb|self.zs.| Moreover, I stored
the input in \verb|self.act| in the beginning to start using them in
the first hidden layer.

Continuing, we create arrays for each layer for preactivation
and activation. Then, we will continue by looping through the
neurons within each layer. Our preactivation will start from
the bias, and by following the definition, we will add the products
of weights and the previous activation. For the if-else condition,
I decided to do this because we don't want $\tanh$ for our final
output layer. Then, we will store the layers in \verb|self.act| and
\verb|self.zs| for latter use.

\begin{jupyter}
\begin{lstlisting}[style=jin]
def backprop(self, y):
    self.grad_w = [[[0.0 for _ in range(len(self.ws[i][j]))]
                    for j in range(len(self.ws[i]))]
                   for i in range(len(self.ws))]
    self.grad_b = [[0.0 for _ in range(len(self.bs[i]))]
                   for i in range(len(self.bs))]
    ypred = self.act[-1][0]
    delta = [0] * len(self.ws)
    delta[-1] = [ypred - y]

    # range(start, stop, step)
    for k in range(len(self.ws) - 2, -1, -1):
        delta[k] = []
        for i in range(len(self.ws[k])):
            dlda = 0.0
            for j in range(len(self.ws[k + 1])):
                dlda += self.ws[k + 1][j][i] * delta[k + 1][j]
            delta[k].append(dlda * tanh_der(self.zs[k][i]))

    for i in range(len(self.ws)):
        prev = self.act[i]
        for j in range(len(self.ws[i])):
            self.grad_b[i][j] = delta[i][j]
            for k in range(len(prev)):
                self.grad_w[i][j][k] = delta[i][j] * prev[k]
\end{lstlisting}
\begin{lstlisting}[style=jout]
\end{lstlisting}
\end{jupyter}

Ah the backprop. I think this is the hardest part from our code.
Assuming that you have read the previous section, let's interpret
the code with the corresponding math equations.

Recall although we do not calcuate the loss per individual parameter,
we do have to calculate the gradients. To start, I decided to put
the initial values as $0$. To store them, we can do the same just
as we randomized the initial weights and biases because ultimately,
each gradient will have a matching parameter. Continuing, \verb|ypred|
is the prediction, and that is the first element from the last
activation layer. This will vary when we have numerous outputs,
but since this is a simple regression problem and we only have one
output, this is good. \verb|delta| here is the delta term for our
hidden layers. Let's start by setting them as $0$ at first and
\verb|ypred - y| for the last layer.

Here comes the fun part. In Python, the \verb|range| function can be
used as \verb|range(start, stop,| \verb|step)|. Here, we start from
the last hidden layer and we subtract $2$ due to indices. Then,
we will stop when we reach $-1$ by going backward by $1$. We will
loop through them for our delta terms. Recall the formula that we
derived below.
\[
\delta_i^{[l]}
= \frac{\partial L}{\partial a_i^{[l]}}
    \frac{\partial a_i^{[l]}}{\partial z_i^{[l]}}
= g'\left(z_i^{[l]}\right) \sum_j w_{ji}^{[l]} \delta_j^{[l+1]}
\]
As the formula suggests, we will add the weights and the
corresponding following delta term with the for loop for
$\frac{\partial L}{\partial a_i^{[l]}}$. Then, we will mulitply
it by the derivative of the actiavtion function, which is $\tanh$
for our example. Finally, we can store our delta terms that are
calculuated. Next, we can now use them to calculate the gradients.
\begin{align*}
\frac{\partial L}{\partial w_{ij}^{[l]}}
&= \delta_i^{[l]} a_j^{[l-1]} \\
\frac{\partial L}{\partial b_i^{[l]}}
&= \delta_i^{[l]} \frac{\partial z_i^{[l]}}{\partial b_i^{[l]}}
= \delta_i^{[l]}
\end{align*}
As the formula suggests, for each parameters, will calculate
the gradients again with for the loop. This is it for the backprop!

\begin{jupyter}
\begin{lstlisting}[style=jin]
def update(self, lr):
    for i in range(len(self.ws)):
        for j in range(len(self.ws[i])):
            for k in range(len(self.ws[i][j])):
                self.ws[i][j][k] -= lr * self.grad_w[i][j][k]
            self.bs[i][j] -= lr * self.grad_b[i][j]

def train(self, train_x, train_y, epochs, lr):
    losses = []
    for epoch in range(epochs):
        total_loss = 0

        for x, y in zip(train_x, train_y):
            ypred = self.forward(x)
            total_loss += loss(y, ypred)
            self.backprop(y)
            self.update(lr)

        average_loss = total_loss / len(train_x)
        losses.append(average_loss)
        
        if (epoch + 1) % 100 == 0:
            print(f"Epoch: {epoch+1}/{epochs}, Loss: {average_loss}")
            
    return losses
\end{lstlisting}
\begin{lstlisting}[style=jout]
\end{lstlisting}
\end{jupyter}

Now we need to update. This part is similar with our previous
example, but for each parameter, we will update it by subracting the
product of the corresponding gradient and the learning rate.
The training part is again very similar to our previous example.
We will begin our loss with $0$, the we will predict, accumulate the
loss, backprop, update, and print every hundred epoch.

\begin{jupyter}
\begin{lstlisting}[style=jin]
    def show_params(self):
        for i, (w, b) in enumerate(zip(self.ws, self.bs)):
            layer = "OUTPUT" if i == len(self.ws) - 1 else f"LAYER {i+1}"
            print(layer)
            for j, neuron_w in enumerate(w):
                weights_str = ", ".join(f"{x:.4f}" for x in neuron_w)
                print(f"  Neuron {j+1}:")
                print(f"    Weights: [{weights_str}]")
                print(f"    Bias = {b[j]:.4f}")
            print()

mlp = MLP([2, 8, 8, 1])
\end{lstlisting}
\begin{lstlisting}[style=jout]
\end{lstlisting}
\end{jupyter}

This part is for printing the parameters, and as you can see, it is
pretty straight forward. I ended up not using it because I don't think
it is that effective at this point, but why not just have it.
Now we create an instance for the MLP class with two hidden layers
with eight neurons each.

\begin{jupyter}
\begin{lstlisting}[style=jin]
losses = mlp.train(train_x, train_y, epochs=2000, lr=0.01)
print()

for x, y in zip(train_x[:10], train_y[:10]):
    print(x, y)
    print(f"     ypred: {mlp.forward(x)}")

plt.plot(losses)
plt.xlabel("epoch")
plt.ylabel("loss")
plt.show()

#mlp.show_params()
\end{lstlisting}
\begin{lstlisting}[style=jout]
\end{lstlisting}
\end{jupyter}

We will train with the learning rate of $0.01$ and $2000$ epochs.
We then print the first ten p9redictions and plot the loss.

\begin{jupyter}
\begin{lstlisting}[style=jin]
test_loss = 0
for x, y in zip(test_x, test_y):
    ypred = mlp.forward(x)
    test_loss += loss(y, ypred)

for x, y in zip(test_x[:10], test_y[:10]):
    ypred = mlp.forward(x)
    print(x, y)
    print(f"    ypred: {ypred}")

print("Test Loss:", test_loss / len(test_x))
\end{lstlisting}
\begin{lstlisting}[style=jout]
\end{lstlisting}
\end{jupyter}

This part is for printing the loss and prediciton for testing.
That was it for the explanation! Now let's take a look at the
results.

\subsubsection{Analyzing the Results}

Now that we have looked at the source code, let's analyze the result
to see how our vanilla neural net performs.

\begin{jupyter}
\begin{lstlisting}[style=jin]
losses = mlp.train(train_x, train_y, epochs=2000, lr=0.01)
print()

for x, y in zip(train_x[:10], train_y[:10]):
    print(x, y)
    print(f"     ypred: {mlp.forward(x)}")

plt.plot(losses)
plt.xlabel("epoch")
plt.ylabel("loss")
plt.show()

#mlp.show_params()
\end{lstlisting}
\begin{lstlisting}[style=jout]
Epoch: 100/2000, Loss: 0.0006795237745380497
Epoch: 200/2000, Loss: 0.0002780199905385479
Epoch: 300/2000, Loss: 0.00015780133554953304
Epoch: 400/2000, Loss: 9.996884725589866e-05
Epoch: 500/2000, Loss: 6.544552166502369e-05
Epoch: 600/2000, Loss: 4.373246439583505e-05
Epoch: 700/2000, Loss: 3.032175895325618e-05
Epoch: 800/2000, Loss: 2.2295750007135315e-05
Epoch: 900/2000, Loss: 1.7559747007674843e-05
Epoch: 1000/2000, Loss: 1.4723606896641237e-05
Epoch: 1100/2000, Loss: 1.2951440585334302e-05
Epoch: 1200/2000, Loss: 1.1773900112740833e-05
Epoch: 1300/2000, Loss: 1.0936643742787753e-05
Epoch: 1400/2000, Loss: 1.0302935457567813e-05
Epoch: 1500/2000, Loss: 9.79811628836463e-06
Epoch: 1600/2000, Loss: 9.379957580934361e-06
Epoch: 1700/2000, Loss: 9.023360655331247e-06
Epoch: 1800/2000, Loss: 8.712547584218995e-06
Epoch: 1900/2000, Loss: 8.437042172730164e-06
Epoch: 2000/2000, Loss: 8.189547717119376e-06

[1.0, -1.0] -1.0
     ypred: -0.9914314084230468
[0.8, -1.0] -0.8
     ypred: -0.8026975789984329
[0.8, -0.9] -0.7200000000000001
     ypred: -0.7229659139930072
[0.1, -0.8] -0.08000000000000002
     ypred: -0.07901899113208088
[0.3, -0.3] -0.09
     ypred: -0.08931568886644314
[-0.7, -0.1] 0.06999999999999999
     ypred: 0.07003157242570901
[-0.6, -0.5] 0.3
     ypred: 0.29863205953269367
[0.7, 0.8] 0.5599999999999999
     ypred: 0.5613061391881997
[-0.5, -0.4] 0.2
     ypred: 0.19629911113434118
[0.9, 0.8] 0.7200000000000001
     ypred: 0.7218073084017129
\end{lstlisting}
\vspace{-\baselineskip}
\includegraphics[scale=0.65]{CS0104.26017-3}
\end{jupyter}

I would say this is pretty good! The loss is constantly decreasing
and the prediction seems to be pretty good. Looking at the graph,
the model seemed to have little idea at first, but instantly learned
to see the holistic view as suggested by the sudden drop.

\begin{jupyter}
\begin{lstlisting}[style=jin]
test_loss = 0
for x, y in zip(test_x, test_y):
    ypred = mlp.forward(x)
    test_loss += loss(y, ypred)

for x, y in zip(test_x[:10], test_y[:10]):
    ypred = mlp.forward(x)
    print(x, y)
    print(f"    ypred: {ypred}")

print("Test Loss:", test_loss / len(test_x))
\end{lstlisting}
\begin{lstlisting}[style=jout]
[-0.9, -0.3] 0.27
    ypred: 0.27224246623233583
[0.9, 1.0] 0.9
    ypred: 0.889162612694173
[0.1, 0.5] 0.05
    ypred: 0.050824340568330406
[1.0, 1.0] 1.0
    ypred: 0.9737260427905209
[0.6, -0.8] -0.48
    ypred: -0.4811987339447982
[0.6, -0.3] -0.18
    ypred: -0.17587974783858806
[0.6, 0.4] 0.24
    ypred: 0.23577905969192914
[0.5, -0.8] -0.4
    ypred: -0.4007379739663608
[0.5, 0.2] 0.1
    ypred: 0.0970149401772119
[0.9, -0.2] -0.18000000000000002
    ypred: -0.18071247280108016
Test Loss: 1.1432408362631248e-05
\end{lstlisting}
\end{jupyter}

Our test also seems pretty good. The loss seems small because our
validation set uses small numbers, but we could see that the
prediction somewhat matches the real values!

Now looking back at our model, this is again good for the limited
range as suggested by the result above. However, there is a reason
for that. Recall that we shuffled the dataset and seletcted $80\%$
for training and $20\%$ for testing there. Therefore, for the limited
range that the model is trained in, the model did pretty good job.
However, when I first tried making dataset by splitting the training
and validation as first $80\%$ and latter $20\%$, our model did
poorly. This is because the model learned the left side well, but
it does poorly in predicting the right side since it has no idea
what's going on there with absolutely no training example. However,
the good news is that our neural net is scalable unlike the previous
example! By increasing the range for datasets, and could see whether
the model underfits or overfits.

We now covered two models from scratch! Now that we know what
happens in deep neural nets, let's build our own language models!
\end{document}


